\subsection{Exhaustive induction profiles at the quartic frontier}
\label{subsec:round4-quartic-profiles}
The two smallest residual tensors admit a substantially smaller finite
search space than the unrestricted central systems suggest.  We enumerate
complete induction multiplicity profiles before searching for additional
twists or fusion coefficients.  A profile specifies the multiplicity of
every possible restriction column; it does not specify a fusion category.

Fix a low induction frame $L$ and let $v_1,\ldots,v_t$ be the remaining
admissible high restriction columns.  If $m_j$ is the number of central
simples with column $v_j$, then
\[
 \sum_{j=1}^t m_jv_jv_j^{\mathsf T}=G,
 \qquad m_j\in\ZZ_{\geq0},
 \qquad
 m_j\leq \min_{i:(v_j)_i>0}
       \left\lfloor\frac{G_{ii}}{(v_j)_i^2}\right\rfloor,
\]
where $G$ is the induction Gram matrix after subtracting $LL^{\mathsf T}$.
The joint character spectrum supplies additional linear equations equating
Galois multiplicities.  All these equations have integer coefficients.

\begin{proposition}[Exact bounded affine enumeration]
\label{prop:round4-affine-profile-enumeration}
For an integer system $Am=b$ with prescribed finite bounds
$0\leq m_j\leq u_j$, row reduction over $\QQ$, followed by the bounded
search below, enumerates every integral solution exactly once.
\end{proposition}
\begin{proof}
Choose the nonpivot coordinates as free variables.  Each pivot coordinate
is an affine rational function of them.  Clear its denominator separately.
Enumerate the free coordinates in any fixed order within their given
integer bounds.  At a partial assignment, the unused coordinates give an
explicit lower and upper bound for each pivot numerator: for a coefficient
$a$ and remaining variable in $[0,u]$, its contribution belongs to
$[\min(0,au),\max(0,au)]$.  Discard a node only if one such interval misses
the allowed interval for that numerator.  At a leaf check divisibility by
each cleared denominator and the original equations.  Every integral
solution follows a unique unpruned path; every retained leaf satisfies the
original system and all bounds.  A node or time budget interrupts this
argument and is therefore recorded as incomplete, never as an exclusion.
\end{proof}

There is a useful inexpensive sieve after enumeration.  Write $D$ for the
global dimension of the original pseudo-unitary category and $d_Z$ for a
central dimension.  If $A_1,\ldots,A_s$ are the nonunit low central simples,
put $\lambda_a(Z)=S_{A_a,Z}/S_{\mathbf1,Z}$, with the normalized modular
matrix $S$.  The character and restriction data determine these numbers.
For every word $w=(a_1,\ldots,a_k)$, simultaneous diagonalization of the
fusion matrices gives
\begin{equation}
\label{eq:round4-word-multiplicity}
 \dim\operatorname{Hom}(\mathbf1,A_{a_1}\otimes\cdots\otimes A_{a_k})
 =\sum_Z\frac{d_Z^2}{D^2}\prod_{j=1}^k\lambda_{a_j}(Z)
 \in\ZZ_{\geq0}.
\end{equation}
No duality assumption on the $A_a$ is needed.  In degree three these are
precisely the low fusion multiplicities after the appropriate dual label
is inserted.  Their restrictions must leave a nonnegative remainder in
every low product.  Both tests are exact linear evaluations of a fixed
profile.  We subsequently check all nonunit low words through degree six.

\begin{proposition}[Finite quartic frontier]
\label{prop:round4-quartic-frontier}
For each of the two surviving low-character labellings, the following
numbers of complete induction profiles remain necessary:
\[
\begin{array}{c|r|r|r|r|l}
\text{tensor}&\text{Gram/Galois}&\text{degree three}&
 \text{words through six}&\text{certified final}&\text{center ranks}\\\hline
R6\_4\_22&93&3&2&2&30,36\\
R6\_4\_91&524&6&6&2&36,42
\end{array}
\]
Any categorification of $R6\_4\_22$ has exactly one central simple whose
restriction is its nontrivial invertible object.  For $R6\_4\_91$, both
simples $g,\sigma$ of its Ising subcategory have exactly one central lift.
These lifts generate an Ising fusion subcategory of its center.
\end{proposition}
\begin{proof}
The general sphericalization and Galois-to-pseudo-unitary reductions are
replayed in the accompanying records, so it suffices to impose the
pseudo-unitary necessary conditions.  All low frames and all second
indicator choices are covered.  The previous exact modular trace rays
remove the other frames and indicator branches; their rational dual
inequalities are replayed here, rather than inferred from numerical
solver status.

For the remaining low frame, Proposition~\ref{prop:round4-affine-profile-enumeration}
visits respectively $20791$ and $4295$ exact search nodes, yielding the
first column of counts.  Equation~\eqref{eq:round4-word-multiplicity} and
restriction products give the next two columns.  One of the three
intermediate profiles for $R6\_4\_22$ has multiplicity $9/2$ for the word
$(1,1,1,1)$ and is discarded.

We then impose the pointed or Ising subcategory projection, the coupled
Galois--twist equations, and the low fusion and restriction equations.
For $R6\_4\_91$, four of the six profiles fail an exact congruence modulo
$2$, on both labellings.  Each certificate is a specified linear
combination of integer-count equations: its coefficient vector vanishes
modulo $2$, while its right side is $1$.  The two remaining profiles have
exact nonnegative witnesses with the required integral multiplicities.
For $R6\_4\_22$, the rank-$36$ profile has such a witness.  The rank-$30$
profile is retained because the bounded search has not supplied an exact
rejection; no feasibility assertion is made for that profile.

Finally, inspection of the exhaustively verified profiles gives
multiplicity one for each stated pure restriction column.  In the Ising
case write the lifts as $G,\Sigma$.  The unique pure unit column is the
central unit.  Uniqueness makes $\Sigma$ self-dual, so rigidity supplies
the unit summand of $\Sigma^2$; the remaining forgotten object is $g$.
Duality and restriction therefore give
$G^2=\mathbf1$, $G\Sigma=\Sigma$ and
$\Sigma^2=\mathbf1+G$, proving the last assertion.  These statements descend
from a pseudo-unitary Galois conjugate because the fusion coefficients and
restriction multiplicities are unchanged.
\end{proof}

The canonical records are
\filename{R6_4_22_frontier_certificate.json} and
\filename{R6_4_91_frontier_certificate.json}, in
\filename{research/round4/quartic_front/}.  The modules
\filename{round4_quartic_profiles.py} and
\filename{round4_quartic_twists.py} reconstruct the finite domains and check
all equalities, congruences, rays and retained witnesses without a
numerical optimizer.  These finite reductions do not construct complete
central fusion data, modular data, or categorifications.  In particular,
both census tensors remain unresolved.
